% !TEX program = lualatex
% Auto-generated from YAML (v3) — 10: データサイエンス・AI（3）：データ分析と仮説検証サイクル

\documentclass[aspectratio=43,professionalfonts,handout]{beamer}
\usepackage{beamer_template}

\newcommand{\LectureNum}{10}
\newcommand{\LectureTitle}{データサイエンス・AI（3）：データ分析と仮説検証サイクル}
\newcommand{\CourseName}{情報リテラシー}
\newcommand{\TextPages}{pp.126-129}
\newcommand{\NextTopic}{データサイエンス・AI（4）：機械学習の基礎}
\newcommand{\NextPages}{-}

\begin{document}
% --- Slide 1: title ---

\begin{frame}{\CourseName\quad 第\LectureNum 回「\LectureTitle 」}
  {\small \TermName\quad \TeacherName\quad ／\quad 教科書 \TextPages}

  \vfill

  \begin{block}{今日の流れ}
    \begin{enumerate}
      \item 仮説検証サイクルとPDCA
      \item 散布図と相関係数
      \item 線形回帰と予測
      \item 訓練データとテストデータ
      \item 相関関係と因果関係・擬似相関
    \end{enumerate}
  \end{block}
\end{frame}

% --- 目標 ---

\begin{frame}{今日の目標}
  \begin{block}{この授業が終わったらできること}
    \begin{itemize}
      \item データ分析の進め方および仮説検証サイクルについて説明できる
    \end{itemize}
  \end{block}
\end{frame}

% --- Slide 2: twocol ---

\begin{frame}{仮説検証サイクルとPDCA}
  \begin{columns}[T]
    \begin{column}{0.48\textwidth}
      \begin{block}{仮説検証サイクル（科学）}
        \small
        \begin{itemize}
          \item ① 仮説：「○○なら△△になるはず」と予想
          \item ② 検証：データで確かめる
          \item ③ 分析：結果を詳しく調べる
          \item ④ 結論：仮説は正しい？修正が必要？
          \item → 新仮説へ（繰り返し改善）
        \end{itemize}
      \end{block}
    \end{column}
    \begin{column}{0.48\textwidth}
      \begin{exampleblock}{PDCA（ビジネス）}
        \small
        \begin{itemize}
          \item Plan（計画）
          \item Do（実行）
          \item Check（確認）
          \item Act（改善）
          \item 共通点：一度で完璧を目指さず，繰り返し改善する
        \end{itemize}
      \end{exampleblock}
    \end{column}
  \end{columns}
\end{frame}

% --- Slide 3: points ---

\begin{frame}{今日の仮説と検証方法}
  \begin{block}{ポイント}
    \begin{itemize}
      \item 仮説：「打率が高い選手ほど，本塁打も多い」
      \item データ：野球選手40人分の成績データ（架空）。セ・リーグ20人，パ・リーグ20人
      \item 項目：選手名・リーグ・打率・本塁打・打点・安打・出塁率
      \item 検証方法1：散布図で視覚的に確認
      \item 検証方法2：相関係数で数値的に確認
      \item 検証方法3：線形回帰で予測モデルを作成
    \end{itemize}
  \end{block}
\end{frame}

% --- Slide 4: terms ---

\begin{frame}{散布図と相関係数}
  \begin{description}[labelwidth=6em]
    \item[\kw{散布図}]
      2つのデータの関係を見るグラフ。横軸と縦軸にそれぞれのデータを取り，点をプロットする
    \item[\kw{相関係数}]
      相関の強さを数値化（-1〜+1）。今日のデータでは約0.7前後 → 比較的高い正の相関
    \item[\kw{注意}]
      相関 ≠ 因果。相関が高くても「打率が本塁打を増やしている」とは断定できない。体格・スイング・戦術など他の要因が関係する可能性がある
  \end{description}
\end{frame}

% --- Slide 5: points ---

\begin{frame}{線形回帰と予測}
  \begin{block}{ポイント}
    \begin{itemize}
      \item 線形回帰：散布図に「最も当てはまる直線（回帰直線）」を引く方法
      \item 数式：y = ax + b（x：説明変数（打率），y：目的変数（本塁打），a：傾き，b：切片）
      \item 決定係数（R²）：今回のデータでは約0.5〜0.6 → モデルが説明できるのは約半分程度
      \item 予測例：打率.300の選手 → 約15〜20本程度／打率.250の選手 → 約10〜15本程度
      \item 外れ値：打率低いのに本塁打多い選手，打率高いのに本塁打少ない選手 → 例外の発見につながる
    \end{itemize}
  \end{block}
\end{frame}

% --- Slide 6: twocol ---

\begin{frame}{訓練データとテストデータ}
  \begin{columns}[T]
    \begin{column}{0.48\textwidth}
      \begin{block}{なぜ分けるのか}
        \small
        \begin{itemize}
          \item モデルが未知のデータでも予測できるか確認するため
          \item 訓練データだけで評価すると，暗記しているだけの可能性がある
          \item 汎化性能（未知データへの対応力）を確認することが重要
        \end{itemize}
      \end{block}
    \end{column}
    \begin{column}{0.48\textwidth}
      \begin{exampleblock}{分け方}
        \small
        \begin{itemize}
          \item 訓練データ（60\%）：モデルの学習に使う。このデータでパラメータを決める
          \item テストデータ（40\%）：モデルの評価に使う。未知のデータとして扱う
          \item 流れ：訓練データで学習 → テストデータで評価 → 汎化性能を確認
        \end{itemize}
      \end{exampleblock}
    \end{column}
  \end{columns}
\end{frame}

% --- Slide 7: twocol ---

\begin{frame}{相関関係と因果関係・擬似相関}
  \begin{columns}[T]
    \begin{column}{0.48\textwidth}
      \begin{block}{相関関係と因果関係}
        \small
        \begin{itemize}
          \item 相関関係：一方が増加すると他方が増加または減少する関係
          \item 因果関係：一方が原因となり他方が結果となる関係
          \item 因果関係の3規準：①共変性（相関あり）②時間的順序③交絡因子の排除
          \item 相関があっても因果関係があるとは限らない
        \end{itemize}
      \end{block}
    \end{column}
    \begin{column}{0.48\textwidth}
      \begin{exampleblock}{擬似相関（見かけの相関）}
        \small
        \begin{itemize}
          \item 例：アイスの売上↑の月は水難事故↑ → 散布図で相関がある
          \item しかしアイスが水難事故を引き起こすわけではない
          \item 交絡因子（気温）が両方に作用しているため
          \item → これを擬似相関という
        \end{itemize}
      \end{exampleblock}
    \end{column}
  \end{columns}
\end{frame}

% --- Slide 8: practice ---

\begin{frame}{Colab実習：野球データの分析}
  {\small 実際にコンピュータで操作してください。}

  \vfill

  \begin{block}{手順}
    \begin{enumerate}
      \item Teamsの課題からリンクをクリック
      \item 「baseball\_data.csv」を自分のマイドライブにコピー
      \item セル1を実行（Googleドライブへのアクセス許可）
      \item セル2〜6を順番に実行
      \item 実行順序：①データ読み込み→②散布図→③相関係数→④train/test分割→⑤線形回帰→⑥予測と可視化
    \end{enumerate}
  \end{block}
\end{frame}

% --- Slide 9: assignment ---

\begin{frame}{課題・次回予告}
  \begin{importantbox}[課題（次回までに提出）]
    WordファイルにスクリーンショットとコメントをまとめてTeamsに提出すること
  \end{importantbox}

  \vfill

  \begin{block}{次回}
    「\NextTopic 」（教科書 \NextPages ）
  \end{block}

  \vfill

  {\small
  質問があれば授業後またはTeamsで受け付けます。}
\end{frame}

% --- チェックリスト ---

\begin{frame}{まとめ・チェックリスト}
  \begin{block}{自己チェック（できたら $\checkmark$ をつけよう）}
    \begin{itemize}
      \item[$\square$] 仮説検証サイクルとPDCAを使ったデータ分析の流れを理解した
      \item[$\square$] 散布図・相関係数・線形回帰による予測の基本を学んだ
      \item[$\square$] 相関関係と因果関係の違い・擬似相関に注意する重要性を把握した
    \end{itemize}
  \end{block}

  \vfill

  {\small 質問があれば授業後またはTeamsで受け付けます。}
\end{frame}

\end{document}
